% Provenance (never printed): sources/book2/BRIEFS.md, g12-08. The business cycle; potential output and the output gap; the
% long-run aggregate model and the neutrality of money; the short-run AD-AS model with sticky prices; demand and supply
% shocks; the return to the long run. Real numbers: figdata/grade-12/08-aggregate-demand-and-supply/ (WDI US real GDP
% 1960-2024 with a Hodrick-Prescott trend, lambda 100; WDI CPI inflation and GDP growth; NBER peaks and troughs as a cited
% table, src/recessions.csv). Ledger: NBER dates and FAQ, Office of the Historian on the 1973-74 embargo, and the g12-07 rows
% on 2007-09. Kestria's curves in the problem are invented.
\chapter{Aggregate Demand and Supply}\label{ch:g12:aggregate-demand-and-supply}

From 1960 to 2024, American output grew about six and a half times, but not smoothly. Nine times in those years output and jobs fell across the economy for months on
end, and each time they recovered. The shortest of these falls lasted two months, in the spring of 2020; the longest lasted eighteen, from December 2007 to June 2009.
% figdata: gdp (3.431 to 22.731 trillion); recessions (9 rows); ledger: g12:nber-2020, g12:nber-2007

The previous chapters explained the long upward path: capital, ideas and institutions. This chapter is about the swings around it. It shows how they are measured, then
gives two models: one for the long run, in which prices have time to adjust, and one for the short run, in which they do not.

\begin{openingchart}{United States: real GDP (solid) and its trend (dashed), trillions of 2015 dollars, 1960--2024, logarithmic scale. Shaded: recessions, from the
peak month to the trough month.}{World Bank, World Development Indicators (NY.GDP.MKTP.KD); trend computed in this book; \omterm{def:g12:aggregate-demand-and-supply:recession}{recession} dates from the NBER Business Cycle
Dating Committee}
\begin{tikzpicture}
\begin{axis}[omtime, width=0.9\linewidth, height=6.0cm, xmin=1960, xmax=2024, ymode=log, ymin=3, ymax=25, log ticks with fixed point,
  ytick={3,5,10,20}, ylabel={trillions of 2015 dollars}]
\omshadeperiods{figdata/grade-12/08-aggregate-demand-and-supply/out/recessions.csv}
\addplot[omDef, thick] table[x=year, y=gdp] {figdata/grade-12/08-aggregate-demand-and-supply/out/gdp.csv};
\addplot[black, dashed] table[x=year, y=trend] {figdata/grade-12/08-aggregate-demand-and-supply/out/gdp.csv};
\end{axis}
\end{tikzpicture}
\end{openingchart}

\begin{recall}
From the grades 6--9 volume: shocks, and the growth of output over time. From \cref{ch:g10:national-accounts}: real GDP. From \cref{ch:g12:inflation-and-money}: the
\omterm{def:g12:inflation-and-money:quantity-equation}{quantity equation} $MV = PY$, and the long-run link between money and prices. From \cref{ch:g12:growth-and-production-function}: output as a function of capital, labour
and technology.
\end{recall}

\section{The facts of the business cycle}\label{sec:g12:aggregate-demand-and-supply:facts}

\begin{definition}[Business cycle]\label{def:g12:aggregate-demand-and-supply:business-cycle}
The \emph{business cycle}\index{business cycle} is the repeated, irregular rise and fall of output, employment and spending around their long-run path.
\end{definition}

\begin{definition}[Recession and expansion]\label{def:g12:aggregate-demand-and-supply:recession}
A \emph{recession}\index{recession} is a period in which economic activity falls significantly across the economy for more than a few months; it runs from a peak to a
trough. An \emph{expansion}\index{expansion}\phantomsection\label{def:g12:aggregate-demand-and-supply:expansion} runs from a trough to the next peak.
\end{definition}

\begin{measure}[How a recession is dated]
% ledger: g12:nber-def
In the United States, a committee of economists at a private research organisation, the NBER, dates peaks and troughs. Its definition is a significant decline in activity,
spread across the economy, lasting more than a few months: depth, diffusion and duration must each be met to some degree. It looks at monthly series: real income less
transfers, two measures of employment, real consumption, real sales of manufacturers and traders, and industrial production. There is no fixed rule, and the committee waits
until the data are clear: it has announced a peak four months after it happened (February 2020) and a trough twenty-one months after (March 1991). A common shortcut,
two quarters of falling GDP, is not its rule.
\end{measure}

\begin{definition}[Potential output and output gap]\label{def:g12:aggregate-demand-and-supply:potential-output}
\emph{Potential output}\index{potential output} is the output the economy produces when its labour and capital are used at normal rates. The \emph{output gap}\index{output
gap}\phantomsection\label{def:g12:aggregate-demand-and-supply:output-gap} is the difference between actual and potential output, as a percentage of potential: negative
in a slump, positive in a boom.
\end{definition}

\omterm{def:g12:aggregate-demand-and-supply:potential-output}{Potential output} cannot be observed; it must be estimated. The opening chart's trend is one estimate: a smooth line drawn through the data by a standard statistical
filter. Its growth slows from 4.3\,\% a year in the 1960s to 2.0\,\% in the 2010s.
% figdata: adas_key (hp_growth_1960s 4.3, hp_growth_2010s 2)

\begin{omchart}{United States: output gap, actual real GDP minus its trend, in per cent of trend, 1960--2024. Shaded: recessions.}{World Bank, World Development Indicators
(NY.GDP.MKTP.KD); trend computed in this book; NBER}
\begin{tikzpicture}
\begin{axis}[omtime, width=0.9\linewidth, height=5.0cm, xmin=1959, xmax=2025, ymin=-6, ymax=4, ytick={-6,-4,...,4}, ylabel={per cent of trend},
  ybar, /pgf/bar width=2.4pt]
\omshadeperiods{figdata/grade-12/08-aggregate-demand-and-supply/out/recessions.csv}
\addplot[fill=omDef!70, draw=none] table[x=year, y=gap] {figdata/grade-12/08-aggregate-demand-and-supply/out/gap.csv};
\end{axis}
\end{tikzpicture}
\end{omchart}
% figdata: adas_key (gap1975 -3.8, gap1982 -5.2, gap2007 2.9, gap2009 -2.9, gap2020 -3.3)

Every \omterm{def:g12:aggregate-demand-and-supply:recession}{recession} on the chart is followed by a negative gap, the deepest $-5.2\,\%$ in 1982; booms, such as the late 1990s or 2007, open positive ones. The gap
closes again, sometimes quickly, sometimes over years.

\section{The long run: output set by resources}\label{sec:g12:aggregate-demand-and-supply:long-run}

\begin{definition}[Long-run aggregate supply]\label{def:g12:aggregate-demand-and-supply:lras}
The \emph{long-run aggregate supply}\index{long-run aggregate supply} curve shows the output produced at each price level once all prices and wages have adjusted: a
vertical line at \omterm{def:g12:aggregate-demand-and-supply:potential-output}{potential output}, which depends on labour, capital and technology, not on the price level.
\end{definition}

\begin{definition}[Neutrality of money]\label{def:g12:aggregate-demand-and-supply:neutrality}
The \emph{neutrality of money}\index{neutrality of money} is the proposition that a change in the quantity of money changes prices in the same proportion and leaves output,
employment and relative prices unchanged.
\end{definition}

\begin{model}{The long-run aggregate model}
\assumes Prices and wages are flexible; output is set by labour, capital and technology.
\predicts Output is at potential; a rise in the quantity of money raises the price level in the same proportion and changes nothing real.
\tested Consistent with it: across 123 countries over decades, money growth and inflation move almost one for one (\cref{ch:g12:inflation-and-money}),
  though output and velocity are not held fixed across them.
\breaks Within a year or two: output falls in \omterm{def:g12:aggregate-demand-and-supply:recession}{recessions} while prices barely move, which flexible prices would not allow.
\end{model}

\section{The short run: sticky prices}\label{sec:g12:aggregate-demand-and-supply:short-run}

\begin{definition}[Sticky price]\label{def:g12:aggregate-demand-and-supply:sticky-price}
A \emph{sticky price}\index{sticky price} is a price, or a wage, that changes only at intervals, because contracts fix it, price lists are costly to change, or sellers
wait to see whether a change in demand will last.
\end{definition}

\begin{definition}[Aggregate demand]\label{def:g12:aggregate-demand-and-supply:ad}
\emph{Aggregate demand}\index{aggregate demand} is the total quantity of output that households, firms, the government and foreigners want to buy at each price level.
\end{definition}

The \omterm{def:g12:aggregate-demand-and-supply:ad}{aggregate demand} curve slopes down. With a given quantity of money and velocity, the \omterm{def:g12:inflation-and-money:quantity-equation}{quantity equation} $PY = MV$ says that spending in money is fixed: a higher price
level leaves room for less real output. Higher prices also make domestic goods dearer than foreign ones.

\begin{definition}[Short-run aggregate supply]\label{def:g12:aggregate-demand-and-supply:sras}
The \emph{short-run aggregate supply}\index{short-run aggregate supply} curve shows the output firms produce at each price level while some prices and wages are fixed:
it slopes up, because a higher price level, with wages and other costs fixed for a time, makes more production worth while.
\end{definition}

\begin{definition}[Aggregate demand and supply shocks]\label{def:g12:aggregate-demand-and-supply:ad-shock}
An \emph{aggregate demand shock}\index{aggregate demand shock} is an event that shifts the \omterm{def:g12:aggregate-demand-and-supply:ad}{aggregate demand} curve, such as a fall in household or business spending. An
\emph{aggregate supply shock}\index{aggregate supply shock}\phantomsection\label{def:g12:aggregate-demand-and-supply:as-shock} is an event that shifts the aggregate supply
curves, such as a rise in the price of an input or a bad harvest.
\end{definition}

\section{Demand shocks and supply shocks}\label{sec:g12:aggregate-demand-and-supply:shocks}

\begin{omfigure}
\begin{tikzpicture}
\begin{groupplot}[group style={group size=3 by 1, horizontal sep=0.5cm}, width=0.37\linewidth, height=4.6cm, xmin=70, xmax=125, ymin=60, ymax=140,
  xtick={90,100}, ytick={80,90,100,110}, xlabel={output}, tick label style={font=\scriptsize}, label style={font=\scriptsize}, grid=major, grid style={gray!20}]
\nextgroupplot[title={\scriptsize a demand shock}, ylabel={price level}]
\addplot[gray, thick] coordinates {(100,60) (100,140)};
\addplot[omThm, thick, domain=70:125, samples=2] {x};
\addplot[omDef, thick, domain=70:125, samples=2] {200 - x};
\addplot[omDef, thick, dashed, domain=70:115, samples=2] {180 - x};
\node[font=\tiny, omThm, anchor=east] at (axis cs:110,118) {SRAS};
\node[font=\tiny, omDef, anchor=south west] at (axis cs:114,87) {AD};
\node[font=\tiny, omDef, anchor=west] at (axis cs:106,68) {AD$'$};
\fill (axis cs:100,100) circle (1.4pt) (axis cs:90,90) circle (1.4pt);
\nextgroupplot[title={\scriptsize a supply shock}, yticklabels={}]
\addplot[gray, thick] coordinates {(100,60) (100,140)};
\addplot[omThm, thick, domain=70:125, samples=2] {x};
\addplot[omThm, thick, dashed, domain=70:115, samples=2] {x + 20};
\addplot[omDef, thick, domain=70:125, samples=2] {200 - x};
\node[font=\tiny, omThm, anchor=north west] at (axis cs:112,108) {SRAS};
\node[font=\tiny, omThm, anchor=east] at (axis cs:96,123) {SRAS$'$};
\node[font=\tiny, omDef, anchor=south west] at (axis cs:114,87) {AD};
\fill (axis cs:100,100) circle (1.4pt) (axis cs:90,110) circle (1.4pt);
\nextgroupplot[title={\scriptsize the return to the long run}, yticklabels={}]
\addplot[gray, thick] coordinates {(100,60) (100,140)};
\addplot[omThm, thick, domain=70:125, samples=2] {x};
\addplot[omThm, thick, dashed, domain=80:125, samples=2] {x - 20};
\addplot[omDef, thick, dashed, domain=70:115, samples=2] {180 - x};
\node[font=\tiny, anchor=west] at (axis cs:100,134) {LRAS};
\node[font=\tiny, omThm, anchor=east] at (axis cs:110,118) {SRAS};
\node[font=\tiny, omThm, anchor=north west] at (axis cs:112,90) {SRAS$''$};
\node[font=\tiny, omDef, anchor=west] at (axis cs:106,68) {AD$'$};
\fill (axis cs:90,90) circle (1.4pt) (axis cs:100,80) circle (1.4pt);
\end{groupplot}
\end{tikzpicture}

{\small Left: a fall in demand from AD to AD$'$ lowers output from 100 to 90 and the price level from 100 to 90. Centre: a rise in costs shifts SRAS up by 20; output falls
to 90 and prices rise to 110. Right: after the demand shock, wages and prices fall over time, SRAS shifts down, and output returns to potential at a price level of 80.}
\end{omfigure}

The two shocks have different signatures. A demand shock moves output and prices in the same direction; a supply shock moves them in opposite directions, the combination
called stagflation. The record shows both.

\begin{record}[2008--09: a fall in demand]
% ledger: g12:nber-2007, g12:gr2007; figdata: adas_key (growth2009 -2.6, cpi2009 -0.4, gap2009 -2.9)
After the fall in house prices and the bank failures of \cref{ch:g12:saving-and-investment}, households and firms cut their spending. The \omterm{def:g12:aggregate-demand-and-supply:recession}{recession} ran from December 2007
to June 2009. Real GDP fell by 2.6\,\% over the year 2009, and consumer prices fell by 0.4\,\%: output and prices moved down together, as a fall in demand predicts. The
\omterm{def:g12:aggregate-demand-and-supply:potential-output}{output gap} reached $-2.9\,\%$.
\end{record}

\begin{record}[1973--74: a supply shock]
% ledger: g12:embargo1973, g12:nber-1973; figdata: adas_key (cpi1973 6.2, cpi1974 11.1, growth1974 -0.5, growth1975 -0.2)
During the Arab--Israeli war of 1973, Arab members of OPEC placed an embargo on oil sales to the United States, lifted in March 1974; the price of a barrel of oil first
doubled, then quadrupled. Oil is an input to almost everything. American consumer price inflation rose from 6.2\,\% in 1973 to 11.1\,\% in 1974 while real GDP fell by
0.5\,\% in 1974 and by 0.2\,\% in 1975; the \omterm{def:g12:aggregate-demand-and-supply:recession}{recession} ran from November 1973 to March 1975. Prices up, output down: the signature of a supply shock.
\end{record}

In 2020 the two kinds of shock arrived together: lockdowns stopped production and spending at once. The \omterm{def:g12:aggregate-demand-and-supply:recession}{recession} lasted only two months, February to April, but real GDP
fell by 2.1\,\% over the year. In 2021 and 2022 output recovered fast while inflation rose to 4.7\,\% and then 8.0\,\%: strong spending met supplies that had not
recovered.
% ledger: g12:nber-2020; figdata: adas_key (growth2020 -2.1, cpi2021 4.7, cpi2022 8)

\section{The return to the long run}\label{sec:g12:aggregate-demand-and-supply:return}

\omterm{def:g12:aggregate-demand-and-supply:sticky-price}{Sticky prices} do not stay stuck. When output is below potential, unemployed workers and idle machines push wages and costs down; the short-run supply curve shifts down,
the price level falls, and output rises back towards potential (the right panel of the figure). When output is above potential, the shift goes the other way. In the
long run, a demand shock changes only the price level, as in the long-run model; in the short run, it changes output.

\begin{model}{The short-run AD--AS model}
\assumes Some prices and wages are sticky; \omterm{def:g12:aggregate-demand-and-supply:ad}{aggregate demand} slopes down, short-run supply slopes up, and long-run supply is vertical at \omterm{def:g12:aggregate-demand-and-supply:potential-output}{potential output}.
\predicts A fall in demand lowers output and prices; a supply shock raises prices and lowers output; with time, output returns to potential.
\tested Consistent with it: 2008--09 (output and prices down together) and 1973--74 (prices up, output down) fit its signatures; in 2021 a demand recovery met a supply shock, and both
  inflation and output rose.
\breaks It does not say how long prices take to adjust, which is the heart of the matter; and expectations of inflation shift the curves themselves, as a later chapter
  shows.
\end{model}

\section{Exercises}

\begin{exercise}[$\star$]\label{exo:g12:aggregate-demand-and-supply:1}
From the opening chart, name the \omterm{def:g12:aggregate-demand-and-supply:recession}{recessions} since 2000, with their years.
\end{exercise}

\begin{exercise}[$\star$]\label{exo:g12:aggregate-demand-and-supply:2}
From the dates in this chapter, how long, in months, were the \omterm{def:g12:aggregate-demand-and-supply:recession}{recessions} of 1973--75 and 2007--09, counted from the peak to the trough?
\end{exercise}

\begin{exercise}[$\star$]\label{exo:g12:aggregate-demand-and-supply:3}
\omterm{def:g12:aggregate-demand-and-supply:potential-output}{Potential output} is 100 and actual output 95. Find the \omterm{def:g12:aggregate-demand-and-supply:potential-output}{output gap}.
\end{exercise}

\begin{exercise}[$\star$]\label{exo:g12:aggregate-demand-and-supply:4}
\omterm{def:g12:aggregate-demand-and-supply:potential-output}{Potential output} is 10 trillion and actual output 10.6 trillion. Find the \omterm{def:g12:aggregate-demand-and-supply:potential-output}{output gap}, and say whether the economy is in a slump or a boom.
\end{exercise}

\begin{exercise}[$\star$]\label{exo:g12:aggregate-demand-and-supply:5}
Classify each shock as a demand shock, a supply shock, or both: a stock-market crash that makes households spend less; a harvest failure; a rise in orders from abroad;
a rise in the price of oil; a lockdown that closes shops and factories.
\end{exercise}

\begin{exercise}[$\star\star$]\label{exo:g12:aggregate-demand-and-supply:6}
\omterm{def:g12:aggregate-demand-and-supply:ad}{Aggregate demand} is $Y = 300 - 2P$ and short-run supply $Y = 2P - 100$. Find the price level and output.
\end{exercise}

\begin{exercise}[$\star\star$]\label{exo:g12:aggregate-demand-and-supply:7}
In \cref{exo:g12:aggregate-demand-and-supply:6}, demand falls to $Y = 260 - 2P$. Find the new short-run price level and output.
\end{exercise}

\begin{exercise}[$\star\star$]\label{exo:g12:aggregate-demand-and-supply:8}
In \cref{exo:g12:aggregate-demand-and-supply:7}, \omterm{def:g12:aggregate-demand-and-supply:potential-output}{potential output} is 100. Find the \omterm{def:g12:aggregate-demand-and-supply:potential-output}{output gap} in the short run, and the price level and output in the long run.
\end{exercise}

\begin{exercise}[$\star\star$]\label{exo:g12:aggregate-demand-and-supply:9}
The quantity of money rises by 10\,\%. What happens to prices and output in the long-run model? And in the short run, with \omterm{def:g12:aggregate-demand-and-supply:sticky-price}{sticky prices}?
\end{exercise}

\begin{exercise}[$\star\star$]\label{exo:g12:aggregate-demand-and-supply:10}
American real GDP was 3.431 trillion 2015 dollars in 1960 and 22.731 trillion in 2024. Find its average annual growth rate.
\end{exercise}

\begin{exercise}[$\star\star$]\label{exo:g12:aggregate-demand-and-supply:11}
From the output-gap chart, which year since 1960 had the most negative gap? Give the gaps of 1975, 2009 and 2020.
\end{exercise}

\begin{exercise}[$\star\star\star$]\label{exo:g12:aggregate-demand-and-supply:12}
In 2009 American real GDP fell by 2.6\,\% and consumer prices by only 0.4\,\%. What does the long-run model, with flexible prices, predict a fall in demand does to output?
What part of the model does 2009 contradict, and which part does it leave standing?
\end{exercise}

\begin{exercise}[$\star\star\star$]\label{exo:g12:aggregate-demand-and-supply:13}
In 2021 American real GDP grew by 6.2\,\% and inflation was 4.7\,\%. What would a supply shock alone predict? A demand recovery alone? Which shift was larger in its effect on
output, and how can both prices and output rise?
\end{exercise}

\begin{exercise}[$\star\star\star$]\label{exo:g12:aggregate-demand-and-supply:14}
After the demand fall of \cref{exo:g12:aggregate-demand-and-supply:7}, the \omterm{def:g12:aggregate-demand-and-supply:potential-output}{output gap} halves each year as wages and prices adjust. Starting from the short-run equilibrium,
give output and the price level (on AD $Y = 260 - 2P$) for three years.
\end{exercise}

\begin{exercise}[$\star\star\star$]\label{exo:g12:aggregate-demand-and-supply:15}
A straight trend fitted on the logarithm of American GDP from 1960 to 2024 grows by 2.92\,\% a year; measured from it, the gap is $-12.5\,\%$ in 1960, $+10.6\,\%$ in 2000 and
$-8.2\,\%$ in 2024. Why are these not good measures of the \omterm{def:g12:aggregate-demand-and-supply:potential-output}{output gap}? Why is even the smooth trend of the opening chart not the same thing as \omterm{def:g12:aggregate-demand-and-supply:potential-output}{potential output}?
\end{exercise}

\section{Problem: A Drought in Kestria}

\begin{problem}[{Weekend problem --- a supply shock and the return to the long run}]
\label{pb:g12:aggregate-demand-and-supply:1}
In Kestria, output $Y$ is measured as an index, 100 at potential, and $P$ is the price level. \omterm{def:g12:aggregate-demand-and-supply:ad}{Aggregate demand} is $Y = 200 - P$, and short-run supply is $P = Y + z$, where
$z$ is a cost shift, zero in normal times.

\medskip
\noindent\textbf{Part I --- Normal times.}
\begin{enumerate}
\item With $z = 0$, find output and the price level.
\item What is the \omterm{def:g12:aggregate-demand-and-supply:potential-output}{output gap}?
\item Draw, or describe, the long-run supply curve.
\item Why does \omterm{def:g12:aggregate-demand-and-supply:ad}{aggregate demand} slope down?
\item Why does short-run supply slope up?
\end{enumerate}

\medskip
\noindent\textbf{Part II --- The drought.}
\begin{enumerate}[resume]
\item A drought raises the costs of food and of everything made from it: in year 1, $z = 20$. Find output and the price level.
\item Find the \omterm{def:g12:aggregate-demand-and-supply:potential-output}{output gap} in year 1.
\item By how much has the price level risen? Is this the signature of a demand or a supply shock?
\item If the government raised demand to bring output back to 100 at once, what would the price level be?
\item Without such a response, what moves the economy back?
\end{enumerate}

\medskip
\noindent\textbf{Part III --- Five years.}
\begin{enumerate}[resume]
\item Each year, wages and other costs adjust so that the cost shift falls by 40\,\% of what remains: $z_t = 20 \times 0.6^{t-1}$. Find $z$ in years 1 to 5.
\item Show that output is $Y_t = 100 - z_t/2$.
\item Find output and the \omterm{def:g12:aggregate-demand-and-supply:potential-output}{output gap} in each of the five years.
\item Find the price level in each year.
\item Add up the output lost over the five years, in per cent of a year's \omterm{def:g12:aggregate-demand-and-supply:potential-output}{potential output}.
\end{enumerate}

\medskip
\noindent\textbf{Part IV --- Interpretation.}
\begin{enumerate}[resume]
\item If costs adjusted by 80\,\% a year instead of 40\,\%, what would the gap be in year 2?
\item How does the speed of adjustment change the output lost?
\item If instead demand had fallen to $Y = 180 - P$ with $z = 0$, find output and prices in the short run and the long run.
\item Compare the price movements in the two shocks.
\item State the result of the problem: the \omterm{def:g12:aggregate-demand-and-supply:potential-output}{output gap} in each of five years, and the cumulative output lost.
\end{enumerate}
\end{problem}
